\documentclass[10pt,a4paper]{amsart}
\usepackage[margin=19mm]{geometry}
\usepackage{amsmath,amssymb,mathtools}
\usepackage[scr=rsfs,cal=euler]{mathalfa}
\usepackage{xfrac}
\usepackage[unicode,hidelinks,bookmarks=false]{hyperref}
\newcommand{\ie}{i.e.\ }
\newcommand{\cf}{cf.\ }
\newcommand{\resp}{respectively\ }
\newcommand{\Subsection}{\subsection}
\providecommand{\nobreakdash}{\nobreak\hbox{-}}
\setlength{\emergencystretch}{2em}
\multlinegap0pt
\usepackage{multirow}
\usepackage{amssymb}
\usepackage{algorithm}\usepackage{algpseudocode}
\usepackage{booktabs}
\usepackage{mathtools}
\usepackage{enumerate}
\usepackage{caption}
\usepackage{float}
\newtheorem{lemma}{Lemma}
\title{Nonholonomic Robot Parking by Feedback--- Part I: Modular Strict CLF Designs}
\author{Velimir Todorovski, Kwang Hak Kim, Alessandro Astolfi, Miroslav Krstic}
\date{Source: arXiv:2511.15119v3 (CC BY 4.0)}
\begin{document}
\maketitle

\noindent\emph{Source and licence.} Nonholonomic Robot Parking by Feedback--- Part I: Modular Strict CLF Designs. Version arXiv:2511.15119v3 (CC BY 4.0), distributed by arXiv under the Creative Commons Attribution 4.0 International licence, http://creativecommons.org/licenses/by/4.0/, as recorded in the arXiv licence metadata for this exact version and shown on the abstract page https://arxiv.org/abs/2511.15119v3. Full licence notice: \texttt{licenses/arxiv-2511.15119-CC-BY-4.0.txt}.\par\smallskip

\noindent\textit{Editorial scope.} Counted unit: the article's lemma lemma:sinterm\_upperbound (source0.tex lines 2638--2655). The article's complete original proof of it is delivered with it (source0.tex lines 2657--2671, one \texttt{proof} environment of 15 lines). No mathematics is written, repaired or re-derived: the statement and the proof are the article's own lines, in the article's own notation, and no retained line is substituted or re-flowed.  No further source-local context is needed: every cross-reference of every retained line resolves inside the two delivered spans.  Nothing was omitted from any retained span, so the package carries no omission record.  No retained line contains a citation, so the wrapper carries no bibliography and no bibliography entry.  No retained line contains a figure environment, an image-inclusion command or drawing code, so no image is delivered and the package declares no asset dependency.  Editorial formatting only, in addition: an AMS \texttt{amsart} preamble that restates the article's own declarations and macros used by the retained lines, namely the environments \texttt{lemma} on separate counters, together with the standard packages those lines need.  The retained environments print in this document as Lemma 1 in order of appearance, whereas the article prints the same items with the numbering of its own full document.  The complete native source of the article as distributed in the arXiv e-print for this exact version is bundled unchanged as \texttt{sources/189423/source0.tex}.
\begin{lemma}\label{lemma:sinterm_upperbound}
The following hold for all $k\geq 1$ and $x\in\mathbb{R}$:
\begin{eqnarray}
1-k \frac{\sin(2x)}{2x} &\le& k x^2 \label{eq:bound1}
\\ 
% \quad \forall x \label{eq:inequality}
%\end{align}
%holds for $k \ge 1 $.
% \end{lemma}
% \begin{lemma}
% \label{lemma:tan_inequality}
% The inequality 
%\begin{align}
1-  
k \cos\gamma(1+\cos\gamma) &\leq& 2 (1 + k)  \tan^2\frac{\gamma}{2}\,. \label{eq:bound2}
\end{eqnarray} 
%holds for $k \ge 1$.
\end{lemma}

\begin{proof}
% Verifiable by ChatGPT (or manually).
The derivation of \eqref{eq:bound1} is omitted as elementary and \eqref{eq:bound2} follows directly by substituting $\cos\gamma=\frac{1-\tan^2(\gamma/2)}{1+\tan^2(\gamma/2)}$.
% With this we have that 
% \begin{align*}
% &1-k \cos \gamma (1 + \cos \gamma)\nonumber\\
% &= 1- k   \frac{1 - \tan^2(\gamma/2)}{1 +\tan^2(\gamma/2)} \left(1 + \frac{1 - \tan^2(\gamma/2)}{1 + \tan^2(\gamma/2)}  \right) \nonumber\\
% % &=1- k  \frac{1 - \tan^2(\gamma/2)}{1 +\tan^2(\gamma/2)} \left( \frac{2 }{1 + \tan^2(\gamma/2)}  \right) \nonumber\\
% % &= 1-  2 \frac{k - k\tan^2(\gamma/2)}{(1 +\tan^2(\gamma/2) )^2} \nonumber\\
% &=   \frac{ 1 + 2\tan^2(\gamma/2) + \tan^4(\gamma/2) - 2k + 2k\tan^2(\gamma/2)}{(1 +\tan^2(\gamma/2) )^2}\nonumber \\
% &\le \frac{  \tan^2(\gamma/2)(2 + 2k + \tan^2(\gamma/2) )}{(1 +\tan^2(\gamma/2) )^2} \nonumber\\
% % & \le \frac{2(1+k) \tan^2(\gamma/2)(1 + \tan^2(\gamma/2) )}{ 1 +\tan^2(\gamma/2) }\nonumber\\
% &\le 2(1+k)\tan^2\left(\frac{\gamma}{2}\right).
% \end{align*}
\end{proof}

\end{document}
